% SGA 3, Exposé IX, tradução brasileira-portuguesa.
% Escopo: §8 integral, bibliografia final e notas editoriais 49--55.
% Autoridade: Polo--Gille Exp9-8nov09.pdf, pp. locais 30--32.

\setcounter{footnote}{48}
\SGARefTarget{sga3:ix:section:8}
\begingroup
\renewcommand{\thefootnote}{*}
\renewcommand{\theHfootnote}{ix-section-eight-star}
\section*{8. Subgrupos, grupos quociente e extensões de grupos de
tipo multiplicativo sobre um corpo\protect\footnotemark}
\addcontentsline{toc}{section}{8. Subgrupos, grupos quociente e extensões
de grupos de tipo multiplicativo sobre um corpo}
\footnotetext{Acrescentado em julho de 1969. Este número tardio é independente
das Seções \SGARefLink{sga3:ix:section:3}{3} a
\SGARefLink{sga3:ix:section:7}{7}.}
\endgroup
\setcounter{footnote}{48}

\medskip
\noindent\SGARefTarget{sga3:ix:scholium:8:0}\textbf{Escólio 8.0.}
\footnote{Nota do editor: acrescentou-se este escólio.} ---
\textit{Sejam \(k\) um corpo e \(H\) um \(k\)-esquema em grupos afim.
Denotemos por \(k[H]\) sua álgebra afim e por
\[
  X(H)=\operatorname{Hom}_{k\text{-}\mathrm{gr}}
    (H,\mathbf G_{m,k})
\]
o grupo de caracteres de \(H\). Pelo lema de independência dos
caracteres, \(X(H)\) é uma parte linearmente independente de \(k[H]\).
Por conseguinte, \(H\) é diagonalizável se, e somente se, \(X(H)\) gera
\(k[H]\) como \(k\)-espaço vetorial.}

\medskip
\noindent\SGARefTarget{sga3:ix:proposition:8:1}\textbf{Proposição 8.1.} ---
\textit{Seja \(G\) um esquema em grupos diagonalizável (respectivamente,
de tipo multiplicativo) sobre um corpo \(k\).\footnote{Nota do editor:
acrescentaram-se as palavras «sobre um corpo \(k\)», implícitas no
original. Ademais, «grupo algébrico» foi substituído por «esquema em
grupos», pois não se supõe que \(G\) seja de tipo finito.} Então,
para todo subesquema em grupos \(H\) de \(G\), os esquemas \(H\) e
\(G/H\) são diagonalizáveis (respectivamente, de tipo multiplicativo).}

A Definição \SGARefLink{sga3:ix:definition:1:1}{1.1} reduz
imediatamente o enunciado ao caso sem a ressalva entre parênteses.
Uma passagem simples ao limite, utilizando VI\textsubscript{B}
\SGARefLink{prop:sga3-VIB-11.13}{11.13} e
\SGARefLink{sga3:viii:theorem:3:1}{VIII 3.1}, reduz-nos ao caso em que
\(G\) é de tipo finito sobre \(k\).\footnote{Nota do editor: seria
preciso detalhar essa passagem ao limite? Utilizar
\SGARefLink{sga3:ix:scholium:8:0}{8.0} e também a igualdade
\[
  G/H=\varprojlim_i G_i/H_i
\]
(cf.\ VI\textsubscript{B}, demonstração de
\SGARefLink{thm:sga3-VIB-11.17}{11.17}). Para ver que
\(H=\varprojlim_i(H\cap G_i)\), utiliza-se que \(H\) é fechado em
\(G\)? Utilizando esse fato, pode-se dar uma demonstração direta.}
Por VI\textsubscript{B}
\SGARefLink{thm:sga3-VIB-11.16}{11.16},\footnote{Nota do editor: sem
dúvida seria preciso reescrever VI\textsubscript{B}
\SGARefLink{thm:sga3-VIB-11.16}{11.16} na forma usual, mais
agradável. Em particular, basta um único \(f_i\) a seguir.} pode-se
encontrar uma família finita de elementos não nulos
\SGARefTarget{sga3:ix:eq:8:1:1}
\begin{equation}
  \tag{8.1.1}
  f_i=\sum_m a_{im}m,\qquad a_{im}\in k,
\end{equation}
do anel afim \(k(M)\) de \(G=D_k(M)\), tal que os pontos de \(H\),
com valores em uma \(k\)-álgebra arbitrária \(k'\), são precisamente os
pontos \(g\in G(k')\) para os quais
\SGARefTarget{sga3:ix:eq:8:1:2}
\begin{equation}
  \tag{8.1.2}
  \tau_g f_i=\lambda_i(g)f_i,
  \qquad \lambda_i(g)\in k',
\end{equation}
onde \(\tau_g\) designa a translação por \(g\). Ora,
\SGARefTarget{sga3:ix:eq:8:1:3}
\begin{equation}
  \tag{8.1.3}
  \tau_g f_i=\sum_m a_{im}\chi_m(g)m,
\end{equation}
onde \(\chi_m:G\to\mathbf G_{m,k}\) é o caráter associado a \(m\).
Portanto, \SGARefLink{sga3:ix:eq:8:1:2}{(8.1.2)} equivale a
\SGARefTarget{sga3:ix:eq:8:1:4}
\begin{equation}
  \tag{8.1.4}
  \chi_m(g)=\chi_{m'}(g)
  \qquad\text{se }m,m'\in Z_i,
\end{equation}
onde \(Z_i\) designa o conjunto dos \(m\in M\) tais que
\(a_{im}\neq0\). A mesma relação pode ser escrita
\[
  \chi_{m'-m}(g)=1
  \qquad\text{se }m,m'\in Z_i.
\]

Seja \(N\) o subgrupo de \(M\) gerado por todos os elementos
\(m'-m\), onde \(i\) varia e \(m,m'\in Z_i\). Da definição de
\(D_k(M)\simeq G\), resulta que \(H\) se identifica com \(D_k(M/N)\).
Portanto, \SGARefLink{sga3:viii:theorem:3:1}{VIII 3.1} identifica
\(G/H\) com \(D_k(N)\).
\hfill\(\square\)

\medskip
\noindent\SGARefTarget{sga3:ix:proposition:8:2}\textbf{Proposição 8.2.}
\footnote{Nota do editor: desenvolveram-se tanto o enunciado quanto sua
demonstração.} ---
\textit{Sejam \(k\) um corpo; \(H\) e \(K\), \(k\)-esquemas em grupos de
tipo multiplicativo e de tipo finito; e \(G\), um \(k\)-esquema em
grupos que entram em uma sequência exata
\SGARefTarget{sga3:ix:eq:8:2:0}
\[
  1\longrightarrow H\longrightarrow G\longrightarrow K\longrightarrow 1
\]
(o que implica que \(G\) é de tipo finito sobre \(k\)).}

\SGARefTarget{sga3:ix:proposition:8:2:item:i}\textit{(i) Se \(G\) é
comutativo ou \(K\) é conexo, então \(G\) é de tipo multiplicativo.}

\SGARefTarget{sga3:ix:proposition:8:2:item:ii}\textit{(ii) Se \(K\) e
\(H\) são diagonalizáveis e \(K\) é um toro, então \(G\) é
diagonalizável.}

Para a demonstração de
\SGARefLink{sga3:ix:proposition:8:2:item:i}{(i)}, remete-se o leitor a
\SGARefLink{sga3:xvii:proposition:7:1:1}{XVII 7.1.1}, do qual
\SGARefLink{sga3:ix:proposition:8:2:item:i}{8.2(i)} é um caso
particular. O caso de um corpo é tratado na parte
\SGARefLink{sga3:xvii:proposition:7:1:1:item:i}{(i)} da demonstração
de \SGARefLink{sga3:xvii:proposition:7:1:1}{XVII 7.1.1}, sem utilizar
os resultados dos exposés posteriores.

Para \SGARefLink{sga3:ix:proposition:8:2:item:ii}{(ii)}, suponhamos
agora que \(K\) e \(H\) são diagonalizáveis:
\[
  K\simeq D_k(M)
  \qquad\text{e}\qquad
  H\simeq D_k(N),
\]
e suponhamos que \(G\) é de tipo multiplicativo, como resulta de
\SGARefLink{sga3:ix:proposition:8:2:item:i}{(i)} quando \(K\) é um
toro. Por \SGARefLink{sga3:x:proposition:1:4}{X 1.4}, \(G\) é
isotrivial sobre \(k\); isto é, existe uma extensão separável finita
\(k'/k\) tal que
\[
  G'=G\times_k k'
\]
é diagonalizável. Por conseguinte, \(G'=D_{k'}(E)\) para um grupo
comutativo \(E\), e tem-se uma sequência exata
\SGARefTarget{sga3:ix:eq:8:2:1}
\begin{equation}
  \tag{1}
  0\longrightarrow D_{k'}(N)\longrightarrow D_{k'}(E)
    \longrightarrow D_{k'}(M)\longrightarrow 0.
\end{equation}
Por \SGARefLink{sga3:viii:theorem:3:1}{VIII 3.1} e
\SGARefLink{sga3:ix:theorem:3:2}{3.2}, \(M\) é, pois, um subgrupo de
\(E\), e tem-se uma sequência exata
\SGARefTarget{sga3:ix:eq:8:2:2}
\begin{equation}
  \tag{2}
  0\longrightarrow M\longrightarrow E\longrightarrow N\longrightarrow 0.
\end{equation}

Para uma extensão dada \(E_0\) de \(N\) por \(M\), consideremos o
\(k\)-grupo diagonalizável \(G_0=D_k(E_0)\). O funtor em grupos sobre
\(k\), \(A\), dos automorfismos da extensão
\SGARefTarget{sga3:ix:eq:8:2:3}
\begin{equation}
  \tag{3}
  1\longrightarrow H\longrightarrow G_0\longrightarrow K\longrightarrow 1,
\end{equation}
é o subfuntor em grupos de
\(\operatorname{Aut}_{k\text{-}\mathrm{gr}}(G_0)\) cujos pontos sobre
um \(k\)-pré-esquema \(T\) são os
\[
  \phi\in\operatorname{Aut}_{T\text{-}\mathrm{gr}}(G_T)
\]
que induzem a identidade sobre \(H_T\) e \(K_T\), identifica-se com
\(\operatorname{Hom}_{k\text{-}\mathrm{gr}}(K,H)\). Por
\SGARefLink{sga3:viii:corollary:1:5}{VIII 1.5}, este é o
\(k\)-grupo constante de valor
\[
  L=\operatorname{Hom}_{\mathrm{gr}}(N,M).
\]
% Candidato a erro de fonte: a extensão é formada com G_0,
% mas o PDF imprime G_T no grupo de automorfismos.

Resulta que as extensões \(G\) de \(K\) por \(H\) que se tornam
isomorfas a \SGARefLink{sga3:ix:eq:8:2:3}{(3)} sobre um fecho
separável \(k_s\) de \(k\) classificam-se do mesmo modo que os
\(k\)-torsores, para a topologia étale, sob o grupo constante \(L_k\).
Ponhamos \(\Gamma=\operatorname{Gal}(k_s/k)\). Como \(L\) é um
\(\Gamma\)-módulo trivial, esses torsores são classificados por
\SGARefTarget{sga3:ix:eq:8:2:4}
\[
  H^1_{\mathrm{\acute et}}(k,L_k)
  =H^1(\Gamma,L)
  =\operatorname{Hom}_{\mathrm{gr.\,top.}}(\Gamma,L).
\]

Se, ademais, \(K\) é um toro, então \(M\) é sem torção e, portanto,
\(L\) também. Por conseguinte,
\[
  \operatorname{Hom}_{\mathrm{gr.\,top.}}(\Gamma,L)=0,
\]
e toda extensão de \(K\) por \(H\) já é diagonalizável.

\medskip
\noindent\SGARefTarget{sga3:ix:remark:8:3}\textbf{Observação 8.3.} ---
\textit{Como já se assinalou na demonstração de
\SGARefLink{sga3:ix:proposition:8:2}{8.2},
\SGARefLink{sga3:ix:proposition:8:2:item:i}{a primeira afirmação} é
enunciada e demonstrada, sobre uma base arbitrária, em
\SGARefLink{sga3:xvii:proposition:7:1:1}{XVII 7.1.1}.
\SGARefLink{sga3:ix:proposition:8:2:item:ii}{A segunda afirmação} será
estendida por descida ao caso de um esquema-base inteiro normal ---ou,
mais geralmente, geometricamente unirramoso--- utilizando
\SGARefLink{sga3:x:lemma:5-13}{X 5.13} mais adiante.}

\textit{Por fim, \SGARefLink{sga3:ix:proposition:8:1}{8.1} também
é um corolário, claramente menos
trivial, de \SGARefLink{sga3:ix:theorem:6:8}{6.8}, que também é válido
sobre um esquema-base arbitrário.}\footnote{Nota do editor: note-se,
contudo, que a demonstração de
\SGARefLink{sga3:ix:theorem:6:8}{6.8} utiliza
\SGARefLink{sga3:ix:proposition:8:1}{8.1}.}

\SGARefTarget{sga3:ix:bibliography}
\section*{Bibliografia}
\addcontentsline{toc}{section}{Bibliografia}
\footnote{Nota do editor: referências adicionais citadas neste
Exposé.}

\noindent \SGARefTarget{sga3:ix:bib:bens}[BEns] N. Bourbaki,
\textit{Théorie des ensembles}, capítulos I--IV, Hermann, 1970.
